\chapter{Overall Discussion and Conclusions}
Please see the discussion and conclusions of each sub-project named in Table \ref{tab:letterciruit_def} (A to G). As previously stated, my recommendations for circuit designs to be included in the `mouse' project going forward are designs B, D and F. In the sub-system sections, specific suggestions have been made regarding future development and project work.

While carrying out the project an additional idea for future project work is to investigate the complete removal of the 9 V alkaline batteries and instead power the entire mouse from 5 x NiMH rechargeable AA batteries. In this configuration, the Arduino would be powered from the AA batteries and supply the voltage required for the operational amplifiers within the position feedback sub-system. 

In conclusion, the aims and objectives of this project have been successfully completed.

{\let\clearpage\relax\chapter{Project Review}}
While the results of the sub-projects are mixed, this has not impacted the overall success of the project with three designs recommended that meet the project objectives defined in section \ref{sec:aims&obj}. The designs suggested were reached after various testing methods and are believed to be the most appropriate circuits to facilitate the reduction in batteries and the addition of feedback to the control loop. Observations and testing conducted during the project have deepened the understanding of the mouse limitations and its potential for further investigation and improvement.

I encountered various issues throughout the project, ranging from testing to component malfunction. Each sub-project section discusses the specific testing issues I encountered. However, one recurring issue encountered during testing was the death of 5 Arduino Nanos, a big setback that knocked my confidence and delayed further tests. I was able to fix this problem by rebuilding my circuits on a separate breadboard, suggesting that the issue lay with the breadboard and not my design.

With hindsight, I would focus on one sub-system at a time, as trying to develop three sub-systems in conjunction leads to mistakes being made in circuit prototyping and missing out on further potential improvements. However, this does not override the positives mentioned previously and hence the project should be classed as a success. 

{\let\clearpage\relax\chapter{Acknowledgment}}
\section{Personal Acknowledgment}
I wish to thank my supervisor, Dr. Francis Robinson, for guiding me through this project and putting up with my endless questions and queries. Additionally, I would like to thank the lab technicians, Richmond Afeawo and Chun JW Wong, who helped me during testing and integration. I would also like to thank Alasdair Philips, who advised me on various topics during testing and the write-up. Lastly, I would like to acknowledge my partner, Bee, who has helped keep me motivated and has been my rock during this project.

\section{AI Acknowledgment}
All statements are based of the suggestions seen in \cite{Rowland_2023}.

I acknowledge that ChatGPT was used to identify themes in the notes I took from the sources I found. I then grouped and organised these themes into a logical order myself and used or adapted some of the themes' phrasing in my topic sentences.

I acknowledge that the Pro version of Grammarly was used to correct grammatical and punctuation errors, and to make minor refinements to my draft text. I reviewed the results and take full responsibility for the content of the assignment.



